% eval_protocol.tex -- Section 10: Evaluation Protocol

\section{Evaluation Protocol}
\label{sec:eval-protocol}

This section describes the comprehensive evaluation protocol used to assess
agent behavior across the full Kant game library, including strategies,
tournament structure, formal metrics, and stratified game splitting.

The headline metric is the agent's mean per-round self-payoff per game
($M_U^{(g)}$, \S\ref{sec:formal-metrics}) --- the same quantity the
training reward optimises, reported in the same units, on a per-game
basis. We do not aggregate across games into a single scalar because
games have different payoff scales. Cooperation rate, fairness, Pareto
efficiency, exploitation resistance, and adaptability are reported as
descriptive secondary outcomes so that the shape of payoff acquisition
can be characterised along the dimensions that matter for alignment
transfer.

\subsection{Strategy Library}
\label{sec:strategy-library}

The evaluation protocol employs 17 opponent strategies spanning two categories.

\begin{table}[h]
\centering\small
\caption{Opponent strategy library. Base strategies (11) apply to all matrix games;
game-specific strategies (6) handle economic games with richer action spaces.}
\label{tab:strategies}
\begin{tabular}{@{}llp{5.8cm}@{}}
\toprule
\textbf{Category} & \textbf{Strategy} & \textbf{Logic} \\
\midrule
\multirow{5}{*}{Unconditional}
In both cases cooperation and defection are considered; for cooperation always play actions MATHI_0 and for defection they use actions MATHI_1
& Random & Uniform random over actions \\
& Mixed & Cooperate w.p.\ $0.7$, defect w.p.\ $0.3$ \\
\midrule
\multirow{5}{*}{Reciprocal}
& Tit-for-Tat & Cooperate first; then mirror opponent \\
& Tit-for-Two-Tats & Defect only after 2 consecutive opponent defections \\
& Grudger & Cooperate until first defection, then always defect \\
& Pavlov & Repeat if same action chosen; else switch \\
& Suspicious TFT & Defect first; then mirror opponent \\
& Generous TFT & Like TFT but cooperate 90\% after defection \\
& Adaptive & Cooperate if opponent cooperated $>50\%$ historically \\
\midrule
Economic and Ultimatum Game - for Pot __MATHI 5__ or __MATHI 6__; Trust Game - generous offers of returns as high as __MATHI 8__ or __MATHI 9__; Public Goods Game - contributions of at least __MATHI 10__ or __MATHI
\bottomrule
\end{tabular}
\end{table}

\subsection{Tournament Structure}
\label{sec:tournament-structure}

Evaluation proceeds via round-robin tournament. For each
$(\text{game}, \text{strategy})$ pair, the \texttt{TournamentRunner}
executes $k$ episodes (default $k=3$) and aggregates results:

\begin{enumerate}[nosep]
    \item For each game $g \in \mathcal{G}$ and strategy $s \in \mathcal{S}$:
    \item \quad For each episode $e \in \{1, \ldots, k\}$:
    \item \quad\quad Reset environment with $(g, s)$; play until terminal state
    \item \quad\quad Record per-round actions, payoffs, cooperation flags
    \item \quad Aggregate: total scores, mean cooperation rate over $k$ episodes
    \item Compute six metrics over the full results matrix (Section~\ref{sec:formal-metrics})
\end{enumerate}

\noindent With $114$ games and $11$ base strategies, a full tournament
executes $114 \times 11 \times 3 = 3{,}762$ episodes.

\subsection{Formal Metrics}
\label{sec:formal-metrics}

We define six metrics. $M_U$ is the headline; $M_C, M_E, M_P, M_F, M_A$
are descriptive secondary outcomes (each normalised to $[0, 1]$);
$M_S$ is retained for backwards comparison only.

\paragraph{Mean Self-Payoff ($M_U$).}
For each (game $g$, opponent strategy $s$) pair, let $S_{g,s}$ be the
total payoff the agent received summed over all $k$ tournament episodes
and $R_{g,s}$ the total number of rounds played in those episodes. Then
\[
M_U^{(g,s)} = S_{g,s} / R_{g,s}
\qquad\text{and}\qquad
M_U^{(g)} = \frac{\sum_s S_{g,s}}{\sum_s R_{g,s}}
\]
is the agent's mean per-round self-payoff against $s$ in $g$, and the
per-game aggregate over all evaluation opponents in $g$. Within one
game $M_U^{(g)}$ is in the same units the training reward
``$r = u_i$'' optimises (\S\ref{sec:reward-shaping}). We deliberately do
\emph{not} aggregate $M_U^{(g)}$ across games into a single scalar:
PD payoffs lie in $\{0,1,3,5\}$, Ultimatum in $[0,10]$, Public Goods
in $[0,\sim 30]$, etc., so a cross-game average mixes incommensurable
units. The headline is the per-game vector $\{M_U^{(g)}\}_g$.

\paragraph{Cooperation Rate ($M_C$).}
Mean fraction of cooperative actions across all $(g, s)$ pairs:
$M_C = \frac{1}{|\mathcal{G}||\mathcal{S}|} \sum_{g,s} \bar{c}_{g,s}$,
where $\bar{c}_{g,s}$ is the cooperation rate averaged over $k$ episodes.
Cooperative actions are $\{\texttt{cooperate}, \texttt{stag}, \texttt{dove}\}$
for matrix games; actions at or above the median index for economic games.

\paragraph{Exploitation Resistance ($M_E$).}
Per-game performance against \texttt{always\_defect} relative to
best/worst:
$M_E^{(g)} = \frac{S_{\text{ad}}^{(g)} - S_{\min}^{(g)}}{S_{\max}^{(g)} - S_{\min}^{(g)}}$,
averaged over games. Score of $1.0$ means no payoff loss to exploitation.

\paragraph{Pareto Efficiency ($M_P$).}
Fraction of $(g,s)$ pairs where joint payoff $u_1 + u_2$ equals the
maximum joint payoff observed for that game:
$M_P = \frac{|\{(g,s) : J_{g,s} \geq J_{\max}^{(g)}\}|}{|\mathcal{G}||\mathcal{S}|}$.

\paragraph{Fairness Index ($M_F$).}
Inverted normalized payoff difference:
$M_F^{(g,s)} = 1 - \frac{|u_1 - u_2|}{|u_1| + |u_2|}$, averaged over
all pairs. Score of $1.0$ indicates perfectly equal payoffs.

\paragraph{Adaptability ($M_A$).}
Normalized variance of cooperation rate across strategies within each game:
$M_A^{(g)} = \frac{\min(\text{Var}_s(\bar{c}_{g,s}),\; 0.5)}{0.5}$,
averaged over games. High variance indicates context-dependent behavior.

\paragraph{Strategic Reasoning ($M_S$, deprecated).}
Unweighted composite: $M_S = \frac{1}{5}(M_C + M_E + M_P + M_F + M_A)$.
Retained for backwards comparison with prior versions of this benchmark;
not used as a headline metric.

\subsection{Stratified Game Splitting}
\label{sec:splitting}

To achieve reliable cross-domain evaluations, we propose domain specific stratification using splitting seeds at $42$ to promote repro

\begin{itemize}[nosep]
    \item \textbf{Train set}: $78\%$ of games ($\approx 89$ games)
    \item \textbf{Eval set}: $22\%$ of games ($\approx 25$ games)
    \item \textbf{Domain constraint}: each domain tag has $\geq 20\%$
          representation in the eval set
\end{itemize}

\noindent The domain constraint ensures no single game category (e.g.,
Evaluation lacks altogether any consideration of classical issues (such as market competition). Together with games that are procedurally produced (see Section \ref{sec:procedural}), this probes if strategic reasoning adapts to new game architectures.
